%----------------------------------------------------------------------------------------------------%
%----------------------------------------------------------------------------------------------------%
%%%% %%%%%%                                     NEWCOMMANDS                                         %%%%%%%%%%%%
%----------------------------------------------------------------------------------------------------%
%----------------------------------------------------------------------------------------------------%
%
% Universale newcommands.tex – alle eigenen Befehle, Farben, Boxen, Kopfzeilen, TikZ- und Listing-Stile.
% Neue Befehle werden HIER ergänzt (nicht im Dokument) und unten protokolliert.
% Quelle: https://github.com/mq90/latex-skills-schule  (CC BY 4.0, Martin Quandt)
%
% Änderungsprotokoll
% 2026-11-09  zusammengeführt aus allen Handreichungen: \TikZ, Box »ausgabe«, Stile für
%             Programmablaufpläne (input, operation, condition, output) und UML (umlclass, inherit,
%             compose, aggregate, associate, depend); Listing-Stil »latex« wieder wie in den
%             Handreichungen ohne Zeilennummern, nummerierte Variante als »latexnummeriert«;
%             Umlaute und · € ° → ≤ ≥ ≠ – — „ “ … × ² ³ π in allen Listings erlaubt (\lstset literate); \lehrerblase, \schuelerblase; PAP-Stile mit align=center (Zeilenumbruch mit \\);
%             \ipa{...} für Lautschrift; \koordinatensystem[Schritt]{xmin}{xmax}{ymin}{ymax};
%             TikZ-Stile punkt, punktoffen, baumdiagramm, wkeit
% 2026-11-09  TikZ-Stil baumdiagrammvertikal (Baumdiagramm von oben nach unten)
% 2026-11-09  Formularfelder für ausfüllbare PDFs: \feld, \zeichen, \kreuz, \antwort, einheitliches
%             Aussehen über \DefaultOptionsofText/CheckBox/Radio/PopdownBox (nur wirksam in einer Form-Umgebung)
% 2026-11-11  Akzentfarben akzentblau, akzentrot, akzentgruen, akzentorange, akzentlila (Karten,
%             Naturwissenschaften, Schulorganisation); Box »hinweis« mit frei wählbarem Titel
% 2026-11-11  Umgebung »passend«: verkleinert den Inhalt auf \linewidth, aber nur wenn er breiter ist
% 2026-11-11  \oz{+II}{Cu}: Oxidationszahl über ein Elementsymbol (Chemie, auch innerhalb \ce{... $\oz{0}{Zn}$ ...})
%








%----------------------------------------------------------------------------------------------------%
%%%%%%%%%%%                                   Mathematische Kürzel                                       %%%%%%%%%%
%----------------------------------------------------------------------------------------------------%


%vereinfachte Schreibweise für den Displaystyle
\newcommand{\dis}[1]{\displaystyle{#1}}


%Zahlenbereiche
\newcommand{\N}{\mathbb{N}} %natürliche Zahlen
\newcommand{\Z}{\mathbb{Z}} %ganze Zahlen
\newcommand{\Q}{\mathbb{Q}} %rationale Zahlen
\newcommand{\R}{\mathbb{R}} %reelle Zahlen
\newcommand{\C}{\mathbb{C}} %komplexe Zahlen
\newcommand{\cP}{\mathcal{P}} %Potenzmenge


%Abkürzungen für Analysis
\newcommand{\eps}{\varepsilon} %Kürzel für Epsilon
\newcommand{\limn}{\underset{n\longrightarrow \infty}{\lim}} %Grenzwert für n gegen Unendlich
\newcommand{\limh}{\underset{h\longrightarrow 0}{\lim}} %Grenzwert für h gegen Null
\newcommand{\limxto}[1]{\underset{x\longrightarrow #1}{\lim}} %beliebiger funktionaler Grenzwert x gegen ....
\newcommand{\dx}{\text{\;d} x} %dx Operator für Integrale
\renewcommand{\d}[1]{\text{\;d}#1} %Operator für beliebige Integrationsvariable
\newcommand{\norm}[1]{\left\lVert#1\right\rVert} %\norm{...} für die Normdarstellung ||...||


%Abkürzungen für Analytische Geometrie
\newcommand{\vektor}[3]{\begin{tiny}\begin{pmatrix} \text{\begin{scriptsize}$#1$\end{scriptsize}}\\ \text{\begin{scriptsize}$#2$\end{scriptsize}}\\ \text{\begin{scriptsize}$#3$\end{scriptsize}}\end{pmatrix}\end{tiny}} %dreidimensionaler Vektor
\newcommand{\punktvektor}[2]{\begin{tiny}\begin{pmatrix} \text{\begin{scriptsize}$#1$\end{scriptsize}}\\ \text{\begin{scriptsize}$#2$\end{scriptsize}}\end{pmatrix}\end{tiny}} %zweidimensionaler Vektor


%Abkürzungen für Stochastik
\newcommand{\E}{\mathbb{E}} %Erwartungswert
\newcommand{\V}{\mathbb{V}} %Varianz


%Umgebung für gegeben und gesucht
\newcommand{\gegges}[2]{%
	\begin{multicols}{2}
		\begin{enumerate}
			\item[geg:] #1
			\item[ges:] #2
		\end{enumerate}
	\end{multicols}
}
%----------------------------------------------------------------------------------------------------%









%----------------------------------------------------------------------------------------------------%
%%%%%%%%%%%%                             Verschiedene  hilfreiche Kürzel                                %%%%%%%%%
%----------------------------------------------------------------------------------------------------%


%römische Darstellung einer Zahl \uproman{2021} ergibt MMXXI
\newcommand{\uproman}[1]{\uppercase\expandafter{\romannumeral#1}}


%Automatisch erstellte Lücke wie die Länge des Textes mit \luecke{<Text>}
\newcommand{\luecke}[1]{\hspace{0,1cm}\underline{\hspace{0,25cm}\hphantom{#1}\hspace{0,25cm}}\hspace{0,1cm}}


%fett-kursive Schrift für neu eingeführte Begriffe
\newcommand{\fk}[1]{\textbf{\textit{#1}}}


%Schwierigkeitsstufen für Aufgaben mit \leicht , \mittel , \schwer
\newcommand{\leicht}{\smash{\raisebox{-1pt}{{\color{black}\rule{4pt}{8pt}}\,{\color{lightgray}\rule{4pt}{8pt}}\,{\color{lightgray}\rule{4pt}{8pt}}}}\,}
\newcommand{\mittel}{\smash{\raisebox{-1pt}{{\color{black}\rule{4pt}{8pt}}\,{\color{black}\rule{4pt}{8pt}}\,{\color{lightgray}\rule{4pt}{8pt}}}}\,}
\newcommand{\schwer}{\smash{\raisebox{-1pt}{{\color{black}\rule{4pt}{8pt}}\,{\color{black}\rule{4pt}{8pt}}\,{\color{black}\rule{4pt}{8pt}}}}\,}
%----------------------------------------------------------------------------------------------------%










%----------------------------------------------------------------------------------------------------%
%%%%%%%                                            Verschiedene Farben                                                      %%%%%%%
%----------------------------------------------------------------------------------------------------%


%Formeln in herkömmlichen Farben schreiben
\newcommand{\frot}[1]{\text{\textcolor{red}{$#1$}}} % rote Formel für Lösungen
\newcommand{\fgruen}[1]{\text{\textcolor{green}{$#1$}}} % grüne Formel für Lösungen
\newcommand{\fblau}[1]{\text{\textcolor{blue}{$#1$}}} % blaue Formel für Lösungen
\newcommand{\fmagenta}[1]{\text{\textcolor{magenta}{$#1$}}} % magenta-farbige Formel für Lösungen
\newcommand{\forange}[1]{\text{\textcolor{orange}{$#1$}}} % orangene Formel für Lösungen
\newcommand{\fcyan}[1]{\text{\textcolor{cyan}{$#1$}}} % cyan-farbige Formel für Lösungen
\newcommand{\folive}[1]{\text{\textcolor{olive}{$#1$}}} % olive-farbige Formel für Lösungen


%Custom made colors
\definecolor{plum}{rgb}{0.36078, 0.20784, 0.4}
\definecolor{chameleon}{rgb}{0.30588, 0.60392, 0.023529}
\definecolor{cornflower}{rgb}{0.12549, 0.29020, 0.52941}
\definecolor{scarlet}{rgb}{0.8, 0, 0}
\definecolor{brick}{rgb}{0.64314, 0, 0}
\definecolor{sunrise}{rgb}{0.80784, 0.36078, 0}
\definecolor{lightblue}{rgb}{0.15,0.35,0.75}
\definecolor{carolina}{rgb}{153, 186, 221}

%Akzentfarben für Grafiken, Karten und Diagramme (kräftig, gut unterscheidbar; aufgehellt z.B. akzentblau!18)
\definecolor{akzentblau}{HTML}{2563EB}
\definecolor{akzentrot}{HTML}{DC2626}
\definecolor{akzentgruen}{HTML}{16A34A}
\definecolor{akzentorange}{HTML}{EA580C}
\definecolor{akzentlila}{HTML}{7C3AED}
%----------------------------------------------------------------------------------------------------%









%----------------------------------------------------------------------------------------------------%
%%%%%%%%                            Kopfzeilen für Leistungsüberprüfungen                              %%%%%%%%%
%----------------------------------------------------------------------------------------------------%


%Definition Kopfzeile für Leistungskontrolle mit \lk{Klasse}{Thema}
\newcommand{\lk}[2]{%
	\pagestyle{fancy}
	\fancyhead[L]{%
	\begin{tabular}[b]{l}
		\\Name:\\Punkte:\\Note:
	\end{tabular} 
	}
	\fancyhead[R]{%
	\begin{tabular}[b]{r}
		\\Leistungskontrolle\\Klasse #1\\ \today
	\end{tabular} 
	}
	\fancyhead[C]{%
	\textbf{\large{#2}}\\
	}
	\fancyfoot[C]{}
	\renewcommand{\headrulewidth}{1pt}
	\fancyheadoffset{1cm}
}


%Definition Kopfzeile für Klassenarbeit mit \ka{Klasse}{Thema}
\newcommand{\ka}[2]{%
	\pagestyle{fancy}
	\fancyhead[L]{%
	\begin{tabular}[b]{l}
		\\Name:\\Punkte:\\Note:
	\end{tabular} 
	}
	\fancyhead[R]{%
	\begin{tabular}[b]{r}
		\\Klassenarbeit\\Klasse #1\\ \today
	\end{tabular} 
	}
	\fancyhead[C]{%
	\textbf{\large{#2}}\\
	}
	\fancyfoot[C]{}
	\renewcommand{\headrulewidth}{1pt}
	\fancyheadoffset{1cm}
}


%Definition Kopfzeile für Klausur mit \kl{Klasse}{Thema}
\newcommand{\kl}[2]{%
	\pagestyle{fancy}
	\fancyhead[L]{%
	\begin{tabular}[b]{l}
		\\Name:\\Punkte:\\Note:
	\end{tabular} 
	}
	\fancyhead[R]{%
	\begin{tabular}[b]{r}
		\\Klausur\\Klasse #1\\ \today
	\end{tabular} 
	}
	\fancyhead[C]{%
	\textbf{\large{#2}}\\
	}
	\fancyfoot[C]{}
	\renewcommand{\headrulewidth}{1pt}
	\fancyheadoffset{1cm}
}


%Entfernt den Header für Leistungskontrollen, Klassenarbeiten oder Klausuren
\newcommand{\clearfancyheader}{%
	\fancyhf{}
	\renewcommand{\headrulewidth}{0pt}
	\renewcommand{\headheight}{0pt}
}


%Definition Kopfzeile für Tägliche Übungen mit \tu{Klasse}{Thema}
\newcommand{\tu}[2]{%
	\pagestyle{fancy}
	\fancyhead[L]{%
	\begin{tabular}[b]{l}
		\\Name:\\Punkte:\\Note:
	\end{tabular} 
	}
	\fancyhead[R]{%
	\begin{tabular}[b]{r}
		\\Tägliche Übung\\Klasse #1\\ \today
	\end{tabular} 
	}
	\fancyhead[C]{%
	\textbf{\large{#2}}\\
	}
	\fancyfoot[C]{}
	\renewcommand{\headrulewidth}{1pt}
	\fancyheadoffset{1cm}
}
%----------------------------------------------------------------------------------------------------%








%----------------------------------------------------------------------------------------------------%
%%%%%%%%%                                     Colorboxes für Hervorhebungen                              %%%%%%%%%
%----------------------------------------------------------------------------------------------------%


 %Colorboxes für Lösungen mit \begin{solution}<Text>\end{solution}
\newtcolorbox{solution}{%
        	lower separated=false,%
	colback=gray!5!white,%
	colframe=gray!75!black,%
        	fonttitle=\bfseries,%
	title=Lösung:,%
      	enhanced,%
	breakable,%
        	attach boxed title to top left={xshift=0.5cm,yshift=-2mm}%
}


%Colorboxes für das chronologische Vorgehen 
\newtcolorbox{todo}{%
        	lower separated=false,%
	colback=yellow!5!white,%
	colframe=orange!75!black,%
        	enhanced,%
	breakable%
}


%Colorboxes für Merksätze
\newtcolorbox{merke}{%
        	lower separated=false,%
	colback=red!5!white,%
	colframe=red!75!black,%
        	fonttitle=\bfseries,%
	title=Merke:,%
      	enhanced,%
	breakable,%
        	attach boxed title to top left={xshift=0.5cm,yshift=-2mm}%
}


%Hinweis-Box mit frei wählbarem Titel: \begin{hinweis}[Lizenz] ... \end{hinweis} (Standardtitel: Hinweis)
\newtcolorbox{hinweis}[1][Hinweis]{%
	enhanced,%
	breakable,%
	sharp corners,%
	colback=akzentblau!5,%
	colframe=akzentblau,%
	boxrule=0pt,%
	leftrule=3pt,%
	fonttitle=\bfseries,%
	coltitle=akzentblau!85!black,%
	attach title to upper={\ },%
	title={#1:}%
}


%Colorbox2 für Merksätze 
\newtcolorbox{merke2}{%
        	lower separated=false,%
        	colback=red!40,%
       	title=Merke,%
       	colframe=black,%
        	fonttitle=\bfseries,%
        	colbacktitle=red,%
        	coltitle=black,%
        	enhanced,%
        	breakable,%
        	attach boxed title to top left={xshift=0.5cm,yshift=-2mm},%
        	title=Merke:%
}
%----------------------------------------------------------------------------------------------------%










%----------------------------------------------------------------------------------------------------%
%%%%%%%%%%%%%                                         MC-Fragen                                    %%%%%%%%%%%%%
%----------------------------------------------------------------------------------------------------%

%Definition des Kästchens zum ankreuzen
\newcommand{\casket}{\LARGE{$\square$}}


%Umgebung für horizontale MC-Fragen
\newenvironment{mcquestionh}[1]{%
	#1%
	\begin{center}\begin{tabular}{>{\centering\arraybackslash}p{2.5cm}  >{\centering\arraybackslash}p{2.5cm} >{\centering\arraybackslash}p{2.5cm} >{\centering\arraybackslash}p{2.5cm} >{\centering\arraybackslash}p{2.5cm}}
	\casket & \casket & \casket & \casket & \casket\medskip\\
	}
	%hier wird die fehlende Tabellenzeile angegeben. Aufruf erfolgt durch:
	%\begin{mcquestionh}{<Frage>}
	%   	\answerh{<Antwort>} & \answerh{<Antwort>} & \answerh{<Antwort>} & \answerh{<Antwort>} & \answerh{<Antwort>}
	%\end{mcquestionh}
	{
	\\
	\end{tabular}\end{center}
}


%Antwort für horizontale MC-Fragen mit \answerh{<Antwort>}
\newcommand{\answerh}[1]{%
	\parbox{2.5cm}{\centering{#1}}
}


%Umgebung für vertikale MC-Fragen
\newenvironment{mcquestionv}[1]{%
    	#1%
    	\begin{enumerate}[\raisebox{-3pt}{\casket}]
    	}
    	%hier wird das fehlende \item angegeben werden. Aufruf erfolgt durch:
    	%\begin{mcquestionv}{<Frage>}
    	%   	\answerv{<Antwort>}
    	%  	\answerv{<Antwort>}
    	%   	\answerv{<Antwort>}
    	%   	\answerv{<Antwort>}
    	%   	\answerv{<Antwort>}
    	%\end{mcquestionv}
    	{
    	\end{enumerate}
}


%Antwort für vertikale MC-Frage mit \answerv{<Antwort>}
\newcommand{\answerv}[1]{%
	\item \hspace{0.1cm} #1
}
%----------------------------------------------------------------------------------------------------%









%----------------------------------------------------------------------------------------------------%
%%%%%%%%%%%                                          Tikz Settings                                          %%%%%%%%%%%%
%----------------------------------------------------------------------------------------------------%


%Koordinatensystem im Standardstil der Mathe-Arbeitsblätter (Gitter gray!30, Achsen -latex, Beschriftung \tiny)
%\koordinatensystem[Schritt]{xmin}{xmax}{ymin}{ymax};
%             TikZ-Stile punkt, punktoffen, baumdiagramm, wkeit
% 2026-11-09  TikZ-Stil baumdiagrammvertikal (Baumdiagramm von oben nach unten) – ganze Zahlen, Bereich muss 0 enthalten,
%Striche und Zahlen bei allen Vielfachen von Schritt (Standard 2). Innerhalb von tikzpicture[scale=0.42] verwenden.
\newcommand{\koordinatensystem}[5][2]{%
	\foreach \kosyx in {#2,...,#3} \draw[gray!30, thin] (\kosyx,#4) -- (\kosyx,#5);
	\foreach \kosyy in {#4,...,#5} \draw[gray!30, thin] (#2,\kosyy) -- (#3,\kosyy);
	\draw[-latex, thick] ({#2-0.3},0) -- ({#3+0.5},0) node[right] {$x$};
	\draw[-latex, thick] (0,{#4-0.3}) -- (0,{#5+0.5}) node[above] {$y$};
	\foreach \kosyx in {#2,...,#3} {%
		\pgfmathtruncatemacro{\kosyrest}{mod(abs(\kosyx),#1)}%
		\ifnum\kosyrest=0 \ifnum\kosyx=0 \else
			\draw (\kosyx,0.1) -- (\kosyx,-0.1) node[below, font=\tiny] {$\kosyx$};
		\fi\fi}
	\foreach \kosyy in {#4,...,#5} {%
		\pgfmathtruncatemacro{\kosyrest}{mod(abs(\kosyy),#1)}%
		\ifnum\kosyrest=0 \ifnum\kosyy=0 \else
			\draw (0.1,\kosyy) -- (-0.1,\kosyy) node[left, font=\tiny] {$\kosyy$};
		\fi\fi}
}


%optionale Eigenschaft für ein  Tikzpicture, damit dieses bündig am oberen Rand einer enumerate-Auflistung steht (benötigt tikzlibrary calc!!!)
\tikzset{
  topline/.style={baseline={([yshift=#1]current bounding box.north)}},
  topline/.default=-\ht\strutbox
}


%Punkte in Graphen und Figuren: \node[punkt, label=above:{$P$}] at (2,3) {}; (offen: punktoffen)
%Als Knoten, damit die Größe nicht mit scale=0.42 mitschrumpft (circle (3pt) würde mitskaliert)
\tikzset{
  punkt/.style={circle, fill, inner sep=0pt, minimum size=4pt},
  punktoffen/.style={circle, draw, fill=white, inner sep=0pt, minimum size=4pt}
}


%Baumdiagramme (Stochastik) von links nach rechts: \begin{tikzpicture}[baumdiagramm] \node[punkt]{} child{...}
%Wahrscheinlichkeiten an den Ästen: edge from parent node[wkeit, above] {$0{,}6$} (oberer Ast above, unterer below)
\tikzset{
  baumdiagramm/.style={grow'=right, sloped, level distance=2.6cm,
    level 1/.style={sibling distance=3cm}, level 2/.style={sibling distance=1.5cm},
    level 3/.style={sibling distance=0.75cm}, edge from parent/.style={draw, thin}},
  wkeit/.style={font=\scriptsize, inner sep=2pt}
}


%Baumdiagramme von oben nach unten: \begin{tikzpicture}[baumdiagrammvertikal] \node[punkt]{} child{...}
%erstes child = linker Ast; Wahrscheinlichkeit links mit node[wkeit, left], rechts mit node[wkeit, right]
\tikzset{
  baumdiagrammvertikal/.style={grow=down, level distance=1.6cm,
    level 1/.style={sibling distance=4.8cm}, level 2/.style={sibling distance=2.4cm},
    level 3/.style={sibling distance=1.2cm}, edge from parent/.style={draw, thin}}
}
%----------------------------------------------------------------------------------------------------%









%----------------------------------------------------------------------------------------------------%
%%%%%%%%%%%                                     Handreichungen                                        %%%%%%%%%%%
%----------------------------------------------------------------------------------------------------%


%Schriftzug TikZ im Fließtext (Leerzeichen danach ist eingebaut)
\newcommand{\TikZ}{Ti\textit{k}Z }


%Oxidationszahl über ein Elementsymbol: $\oz{+II}{Cu}$, in \ce{...} zwischen Dollarzeichen: \ce{$\oz{0}{Zn}$ + ...}
\newcommand{\oz}[2]{\overset{\scriptscriptstyle\mathrm{#1}}{\mathrm{\ce{#2}}}}


%Umgebung passend: Inhalt (z.B. eine große tikzpicture) höchstens so breit wie die Zeile setzen.
%Ist er schmaler, bleibt er unverändert. \begin{passend} ... \end{passend} (Inhalt nicht als Argument, # erlaubt)
\newsavebox{\passendbox}
\newenvironment{passend}{\begin{lrbox}{\passendbox}}{\end{lrbox}%
	\ifdim\wd\passendbox>\linewidth\resizebox{\linewidth}{!}{\usebox{\passendbox}}\else\usebox{\passendbox}\fi}


%Box für die Ausgabe unter einem Code-Beispiel: erst lstlisting[style=latex], darunter \begin{ausgabe} ... \end{ausgabe}
\newtcolorbox{ausgabe}{%
	colback=blue!5,%
	colframe=blue!50!black,%
	title=Ausgabe,%
	fonttitle=\bfseries%
}
%----------------------------------------------------------------------------------------------------%









%----------------------------------------------------------------------------------------------------%
%%%%%%%%%%%                                        Lautschrift                                         %%%%%%%%%%%
%----------------------------------------------------------------------------------------------------%


%Lautschrift in eckigen Klammern mit \ipa{"a:b@nt} – schaltet das babel-Kürzel " vorher aus (sonst wird "a zu ä)
\makeatletter
\newcommand{\ipa}{\begingroup\shorthandoff{"}\ipa@aux}
\newcommand{\ipa@aux}[1]{[\textipa{#1}]\endgroup}
\makeatother
%----------------------------------------------------------------------------------------------------%









%----------------------------------------------------------------------------------------------------%
%%%%%%%%%%%                                       Sprechblasen                                         %%%%%%%%%%%
%----------------------------------------------------------------------------------------------------%


%Sprechblase mit Zipfel links unten (Lehrkraft) mit \lehrerblase[TikZ-Optionen]{Text} (benötigt shapes.callouts)
\newcommand{\lehrerblase}[2][]{%
  \tikz[baseline=(B.base)]{%
    \node[rectangle callout, rounded corners=6pt, callout relative pointer={(-13mm,-7mm)},
      draw, thick, fill=blue!6, align=left, inner sep=3.5mm, #1] (B) {#2};}%
}


%Sprechblase mit Zipfel rechts unten (Schülerin/Schüler) mit \schuelerblase[TikZ-Optionen]{Text}
\newcommand{\schuelerblase}[2][]{%
  \tikz[baseline=(B.base)]{%
    \node[rectangle callout, rounded corners=6pt, callout relative pointer={(13mm,-7mm)},
      draw, thick, fill=green!6, align=left, inner sep=3.5mm, #1] (B) {#2};}%
}
%----------------------------------------------------------------------------------------------------%









%----------------------------------------------------------------------------------------------------%
%%%%%%%%%%%                            TikZ-Stile für Programmablaufpläne                            %%%%%%%%%%%
%----------------------------------------------------------------------------------------------------%


%Ein-/Ausgabe als Parallelogramm, Operation als Rechteck, Verzweigung als Raute (benötigt shapes.geometric)
%align=center, damit \\ im Knotentext eine neue Zeile ergibt (ohne align wird \\ stillschweigend verschluckt)
\tikzstyle{input}=[trapezium, trapezium left angle=60, trapezium right angle=120, fill=green!20, draw=green, align=center]
\tikzstyle{operation}=[rectangle, fill=orange!20, draw=orange, align=center]
\tikzstyle{condition}=[diamond, fill=magenta!20, draw=magenta, align=center]
\tikzstyle{output}=[trapezium, trapezium left angle=60, trapezium right angle=120, fill=red!20, draw=red, align=center]
%----------------------------------------------------------------------------------------------------%









%----------------------------------------------------------------------------------------------------%
%%%%%%%%%%%                            TikZ-Stile für UML-Klassendiagramme                           %%%%%%%%%%%
%----------------------------------------------------------------------------------------------------%


%Klasse mit drei Teilen: Name, Attribute, Methoden (benötigt shapes.multipart)
\tikzstyle{umlclass} = [rectangle split, rectangle split parts=3, rectangle split part align={center,left,left}, draw, thick, text width=4cm, minimum height=2cm]


%Beziehungen (benötigt arrows.meta)
\tikzstyle{inherit} = [-{Triangle[open]}, thick]	%Vererbung
\tikzstyle{compose} = [-{Diamond}, thick]	%Komposition
\tikzstyle{aggregate} = [-{Diamond[open]}, thick]	%Aggregation
\tikzstyle{associate} = [-, thick]	%Assoziation
\tikzstyle{depend} = [-, thick, dashed]	%Abhängigkeit
%----------------------------------------------------------------------------------------------------%








%----------------------------------------------------------------------------------------------------%
%%%%%%%%%%%                                         Listing Settings                                       %%%%%%%%%%%%
%----------------------------------------------------------------------------------------------------%


%Umlaute und häufige Sonderzeichen in ALLEN Listings (Python, Java, ...), sonst Fehler »Invalid UTF-8 byte sequence«
\lstset{literate=%
	{ä}{{\"a}}1 {ö}{{\"o}}1 {ü}{{\"u}}1 {Ä}{{\"A}}1 {Ö}{{\"O}}1 {Ü}{{\"U}}1 {ß}{{\ss}}1%
	{·}{{$\cdot$}}1 {€}{{\euro}}1 {°}{{\degree}}1 {→}{{$\rightarrow$}}2 {≤}{{$\leq$}}1 {≥}{{$\geq$}}1 {≠}{{$\neq$}}1%
	{–}{{--}}1 {—}{{---}}1 {„}{{\glqq}}1 {“}{{\grqq}}1 {…}{{\dots}}3 {×}{{$\times$}}1 {²}{{$^2$}}1 {³}{{$^3$}}1 {π}{{$\pi$}}1%
}


%Einstellung der Codeumgebung für LaTeX-Code-Beispiele (Stil der Handreichungen: ohne Zeilennummern)
%Verwendung: \begin{lstlisting}[style=latex] ... \end{lstlisting}, darunter \begin{ausgabe} ... \end{ausgabe}
\lstdefinestyle{latex}{%
  	language=TeX,%
  	frame=single,%
  	backgroundcolor=\color{yellow!20},%
  	rulecolor=\color{lightgray!40},%
  	breaklines=true,%
  	basicstyle=\scriptsize\ttfamily,%
  	keywordstyle=\bfseries\color{purple!40!black},%
  	commentstyle=\itshape\color{green!40!black},%
  	identifierstyle=\color{blue},%
	stringstyle=\color{black},%
	captionpos=b,%
	tabsize=2,%
	showspaces=false,%
	showstringspaces=false,%
  	showtabs=false,%
	deletekeywords={document},%
	morekeywords={begin}%
}


%Variante mit Zeilennummern: \begin{lstlisting}[style=latexnummeriert]
\lstdefinestyle{latexnummeriert}{%
	style=latex,%
	numbers=left,%
	numberstyle=\tiny\color{gray!90!black},%
  	numbersep=-5pt%
}
%----------------------------------------------------------------------------------------------------%






%----------------------------------------------------------------------------------------------------%
%%%%%%%%%%%                          Formularfelder (ausfüllbare PDFs)                          %%%%%%%%%%%
%----------------------------------------------------------------------------------------------------%
% Nur innerhalb von \begin{Form} ... \end{Form} (hyperref) verwendbar. Für Optionsfelder (radio)
% im Dokument zusätzlich \hypersetup{unicode=false}, sonst bekommen sie kaputte Namen.
% Jedes Feld steckt in einer eigenen Box (\raisebox/\hbox): Ohne sie übernimmt pdfTeX die Höhe
% der umgebenden Zeile, in Tabellen werden die Felder dann so hoch wie die ganze Zeile.

%Aussehen aller Felder: blauer Rand, hellblauer Hintergrund, Schrift 11 pt, wird mitgedruckt
\renewcommand{\DefaultOptionsofText}{print, bordercolor={0 0.353 0.612}, backgroundcolor={0.94 0.96 1}, charsize=11pt}
\renewcommand{\DefaultOptionsofCheckBox}{print, bordercolor={0 0.353 0.612}, backgroundcolor={0.94 0.96 1}, charsize=11pt}
\renewcommand{\DefaultOptionsofRadio}{print, notoggletooff, bordercolor={0 0.353 0.612}, backgroundcolor={0.94 0.96 1}, charsize=11pt}
\renewcommand{\DefaultOptionsofPopdownBox}{print, bordercolor={0 0.353 0.612}, backgroundcolor={0.94 0.96 1}, charsize=11pt}
\renewcommand{\DefaultHeightofText}{0.6cm}

%\feld[Breite]{Name} – einzeiliges Textfeld, sitzt auf der Grundlinie (Standardbreite 2,4 cm)
\newcommand{\feld}[2][2.4cm]{\raisebox{-0.17cm}{\TextField[name=#2, width=#1]{}}}

%\zeichen{Name} – Feld für genau ein Zeichen, zentriert (z. B. E/V/A, Buchstabe einer Reihenfolge)
\newcommand{\zeichen}[1]{\raisebox{-0.12cm}{\TextField[name=#1, width=0.8cm, height=0.5cm, maxlen=1, align=1]{}}}

%\kreuz{Name} – Ankreuzkästchen
\newcommand{\kreuz}[1]{\raisebox{-0.08cm}{\CheckBox[name=#1, width=0.42cm, height=0.42cm]{}}}

%\antwort[Höhe]{Name} – mehrzeiliges Antwortfeld über die ganze Zeilenbreite (Standardhöhe 3,2 cm ≈ 4 Zeilen)
\newcommand{\antwort}[2][3.2cm]{\par\vspace{0.15cm}\noindent\hbox{\TextField[name=#2, multiline=true, width=\linewidth, height=#1]{}}}
%----------------------------------------------------------------------------------------------------%
